\chapter{INTRODUCTION}

\section{Research Background}

The rapid paradigm shift toward high-level autonomous systems is fundamentally
transforming global industrial modalities, establishing new benchmarks for
operational efficiency, safety, and productivity across diverse sectors
\citep{Guizzo_2008_threeengineers}. Within this broader context, off-road
autonomous navigation represents an especially critical and challenging
deployment domain, particularly for heavy agricultural machinery tasked with
operating autonomously within unstructured and severe field environments
\citep{Hameed_2021_comprehensivereview}. Unlike highly regulated urban streets
or uniform industrial indoor settings, unstructured natural terrains present
multi-faceted complexities, including irregular surface roughness, unpredictable
slope variations, highly irregular topographies, and dynamic obstacles
\citep{Sofman_2006_improvingrobot}. To guarantee reliable vehicle state
estimation and successful terrain traversal under these erratic conditions, the
integration of robust, terrain-aware path planning architectures represents an
indispensable algorithmic prerequisite \citep{Sedighi_2019_guidedhybrid}.
Consequently, overcoming these geometric terrain hazards stands as a primary
engineering challenge in modern mobile robotics design.

Within the specific domain of precision agriculture, autonomous navigation
frameworks must directly resolve acute demographic, economic, and operational
constraints. The global agricultural sector is currently facing a dual
challenge: an escalating demand for food and crop products driven by population
growth, occurring alongside a severe, long-term contraction in the available
skilled labor pool \citep{Emmi_2014_newtrends}. To mitigate these macro-economic
vulnerabilities, increasing the scale of dependable in-field automation and
mobile robotics platforms has become a structural necessity
\citep{Emmi_2014_newtrends}. Effective path planning algorithms designed for
these agricultural platforms must look past simple binary (passable/impassable)
obstacle avoidance; they must concurrently minimize mechanical energy
dissipation, optimize total travel time, and enforce absolute platform stability
over highly undulating landscapes \citep{Hameed_2021_comprehensivereview}. To
systematically address these multi-objective criteria, this study focuses on the
development, implementation, and validation of a comprehensive 2.5D terrain-
aware path planning framework explicitly engineered to optimize autonomous
ground vehicle routing across challenging off-road and agricultural
topographies.

\section{Introduction to Path Planning for Autonomous Robots}

Path planning represents a foundational architectural pillar of mobile robot
autonomy, enabling automated platforms to execute localized routing maneuvers
deterministically. Mechanistically, this process is tasked with computing an
optimal or near-optimal trajectory from an initial coordinates origin to a
targeted terminal destination while concurrently satisfying explicit kinematic
constraints and avoiding localized environmental obstacles
\citep{SanchezIbanez_2021_pathplanning}. Historically, classical graph-search
and grid-allocation strategies, most notably the Dijkstra and A* algorithms,
were pioneered to resolve these pathfinding problems within planar, two-
dimensional (2D) static environments \citep{Karur_2021_surveypath,
Noto_2000_methodshortest}. These foundational frameworks operate by evaluating
discrete traversal costs between adjacent nodes distributed across a structured
grid network or topological graph representation, typically optimizing the
search space to minimize cumulative path lengths or similar spatial clearance
metrics \citep{Karur_2021_surveypath, Noto_2000_methodshortest}. For example,
the deterministic search behavior of the standard A* algorithm is governed by a
multi-variant evaluation cost function formalized as:

\begin{equation}
    f(n) = g(n) + h(n)
\end{equation}

\noindent where $g(n)$ defines the exact accumulated cost computed from the initial start node to
the current node $n$, and $h(n)$ represents the heuristic estimation of the remaining traversal cost
required to reach the terminal goal state \citep{Karur_2021_surveypath, Noto_2000_methodshortest}. While these planar graph-search
heuristics provide global mathematical optimality within flat workspace boundaries, their
operational efficacy degrades significantly when deployed across non-planar topographies.

Traditional A* path planning paradigms fall short when confronted with the
complex, uneven terrains typical of off-road and agricultural settings because
they inherently disregard the vertical elevation dimension
\citep{Bettencourt_2025_indoor25d}. This geometric omission can lead to severe
navigation failures, such as catastrophic collisions with unmodeled structural
obstacles positioned physically above the robot's primary 2D reference plane, or
the generation of global paths that are physically untraversable due to
excessive, unsafe slope gradients \citep{Bettencourt_2025_indoor25d}. To resolve
these spatial deficiencies without incurring prohibitive processing overheads,
two-and-a-half-dimensional (2.5D) environmental representations have emerged as
a highly practical structural compromise between full 3D volumetric voxel
complexity and 2D planar simplicity. Mathematically, a 2.5D elevation map
correlates horizontal planar coordinates $(x, y)$ with a single, discrete
altitude value $z$ for every coordinate cell, successfully capturing localized
terrain height variations and allowing the system to derive key topographic
features like slope gradients \citep{Bettencourt_2025_indoor25d}. This strategic
abstraction enables pathfinding algorithms to directly factor in critical
traversability criteria and localized step distances, completely bypassing the
extreme computational burdens and data densities associated with full 3D
volumetric representations \citep{Bettencourt_2025_indoor25d}. Transitioning to
these 2.5D mapping models is particularly relevant for autonomous agricultural
vehicles, where executing stable navigation across sloped landscapes and
unmanaged crop fields requires a highly nuanced, topologically aware
understanding of terrain profiles to prevent platform instability
\citep{Hameed_2021_comprehensivereview}.

\section{Problem Statement}

The primary challenge in off-road and agricultural navigation stems from the
dynamic, non-linear complexities of uneven terrain, which cannot be adequately
modeled or predicted by conventional two- dimensional (2D) path planning
paradigms \citep{Bettencourt_2025_indoor25d}. Relying exclusively on 2D planar
configurations risks generating trajectories that are locally optimal in terms
of horizontal distance but physically impossible, highly inefficient, or
dangerous due to severe slope gradients, large elevation changes, or undetected
vertical obstacles \citep{Bettencourt_2025_indoor25d}. For instance, a planar
navigation routine may execute a command leading to a terminal collision if it
fails to evaluate obstacles positioned above its active sensor observation
plane, inducing severe tracking and navigation errors
\citep{Bettencourt_2025_indoor25d}. Furthermore, attempting to retroactively fix
these geometric oversights via standard path post-processing often fails to
resolve fundamental physical incompatibilities, yielding discontinuous,
``choppy'' trajectories or acute turning angles that present severe vehicle
control challenges \citep{Frey_2022_selfsupervisedneartofar}. Consequently,
relying on flat surface assumptions directly jeopardizes vehicle safety and
structural integrity on natural topographies.

Existing off-road path planning methodologies face a strict trade-off between
computational efficiency and environmental fidelity. Traditional two-dimensional
(2D) algorithms—such as standard A*, Dijkstra, and Rapidly-exploring Random
Trees (RRT) achieve high computational velocity by oversimplifying terrain
models. However, because these 2D approaches treat the environment as a flat
plane, they fail to account for elevation changes, leading to unsafe or
physically impossible paths in complex off-road environments
\citep{Overbye_2022_localoffroad}. Conversely, full three-dimensional (3D)
configurations—such as 3D A*, Hybrid A*, and 3D Occupancy Grid maps—incur
prohibitive computational overhead. These 3D methods suffer from an exponential
expansion of the search state space (x, y, z, and orientation), making real-time
execution intractable on resource-constrained embedded hardware. While
widespread modifications to graph-search algorithms successfully optimize
runtime search velocities, minimize raw path lengths, or manage transient
dynamic obstacles \citep{Injarapu_2023_surveyautonomous, Lei_2021_globalpath,
Li_2020_pathplanning}, they fail to systematically integrate multi-dimensional
terrain characteristics into the core cost-evaluation loop for proactive
decision-making. Specifically, within the domain of agricultural robotics, the
direct, real-time evaluation of physical three-dimensional step distances (3DSD)
and localized slope magnitudes (SM) within the internal A* search mechanism
compounded by an adjustable slope weight remains a critical research gap
requiring rigorous investigation. Additionally, standard trajectory smoothing
techniques routinely execute geometric alterations without enforcing compliance
with hard environmental traversability boundaries, such as the vehicle's Maximum
Climbable Slope (MCS). This algorithmic omission results in smoothed paths that
look visually appealing but violate the platform's actual physical climbing
tolerances.

\section{Research Gap}

Despite extensive advancements in autonomous vehicle navigation and graph-search
adaptations, several critical gaps persist regarding the integrated frameworks
required for safe, efficient, and terrain-aware navigation in complex off-road
and agricultural settings. This overarching deficiency in contemporary field
robotics literature is characterized by three distinct technical gaps:

Firstly, current 2.5D path planning algorithms rarely integrate 3D terrain
features directly into their core cost evaluations. While the classical A*
algorithm has seen numerous architectural enhancements designed to optimize
runtime search efficiency, manage binary obstacle avoidance, and navigate
transient dynamic environments \citep{Injarapu_2023_surveyautonomous,
Lei_2021_globalpath, Li_2020_pathplanning}, fewer studies explicitly embed a
comprehensive set of 2.5D terrain features, such as varying slope magnitudes
(SM) and precise 3D step distances (3DSD), directly into its cost function for
agricultural applications. Existing 2.5D navigation methods predominantly focus
on discrete obstacle representation and structural layout mapping rather than
nuanced traversability costs derived directly from continuous terrain geometry.
Furthermore, contemporary frameworks frequently utilize generic or static cost
indices that fail to dynamically adapt to the highly variable topographies and
specific operational challenges encountered by agricultural robots
\citep{Bettencourt_2025_indoor25d}. Consequently, the development of a unified
multi-terrain cost function that proactively guides the graph search based on
these physical terrain properties and an adjustable slope weight ($\lambda$)
represents an unaddressed and highly necessary contribution.

Secondly, existing path-smoothing techniques do not enforce physical climbing
limits. While smoothing creates visually clean trajectories, standard 2D and 3D
methods ignore vehicle boundaries like maximum climbable slope—often turning a
valid path into an unnavigable hazard. Standard path optimization protocols
widely implement geometric refinement methods, such as B-splines and Bezier
curves, to elevate path aesthetics, eliminate acute turning variations, and
improve vehicle drivability \citep{Han_2019_agvpath, Zhang_2021_mobilerobot,
Frey_2022_selfsupervisedneartofar}. However, a significant deficiency exists in
consistently coupling these continuous geometric smoothing processes with a
rigorous validation of terrain- climbability thresholds (e.g., Maximum Climbable
Slope) and other localized boundary constraints. Because typical smoothing
algorithms focus exclusively on minimizing curvature derivatives, they risk
generating kinematically fluid trajectories that, while visually appealing,
inadvertently route the vehicle through steep sections that violate the actual
physical traversability limits of the platform, leading to unsafe or impossible
execution profiles \citep{Frey_2022_selfsupervisedneartofar}. Currently, the
explicit coupling of B-spline refinement with discrete sampling checks for
boundary compliance, vegetation exposure, and maximum slope thresholds is not
extensively addressed in the literature for 2.5D terrain-aware planning in
agricultural environments. Proactively embedding these physical threshold checks
immediately following or during the smoothing loop is therefore essential to
prevent vehicle instability.

Finally, existing research lacks thorough experimental evaluations across
diverse terrain types. Current frameworks often rely on oversimplified geometric
maps and fail to provide direct, side-by-side comparisons between terrain-aware
2.5D algorithms and standard 2D baselines. For example, comprehensive reviews by
\citet{Hameed_2021_comprehensivereview} highlight how path planning evaluations
routinely overlook complex elevation changes, while studies like
\citet{Song_2019_smoothedalgorithm} focus primarily on 2D path smoothing without
benchmarking against 2.5D terrain-informed cost functions. Consequently,
algorithms remain untested across varied topographic conditions, such as flat
plains, rolling hills, and steep slopes. Furthermore, studies rarely evaluate a
complete set of performance metrics—such as true 3D path length, execution time,
and energy consumption—at the same time. A rigorous comparison across different
slope weights ($\lambda$) and terrain profiles is essential to prove the real-
world value and computational efficiency of 2.5D path planning.

\section{Research Objectives}

The primary objective of this research is to design, implement, and validate a
comprehensive, safety-conscious navigation architecture tailored for autonomous
ground vehicles (AGVs) operating across unstructured environments. To achieve
this macro-level goal, the study is systematically segmented into three primary
technical objectives:

\begin{enumerate}
    \item To develop a 2.5D path planning algorithm that integrates elevation profiles and localized slope constraints into the cost-evaluation loop for multi-terrain environments.

\item To implement a B-spline trajectory smoothing technique that refines the
    planned path while strictly enforcing maximum climbable slope limits.

    \item To validate and benchmark the performance of the proposed algorithm against standard baseline planners across diverse terrain types using 3D path length, path error, execution time, and energy consumption.
\end{enumerate}

\section{Research Questions}

To systematically achieve the stated research objectives and address the
identified gaps in field robotics literature, this study is guided by three core
research questions that isolate the algorithmic, kinematic, and empirical
dimensions of the navigation framework:

\begin{enumerate}
    \item How can three-dimensional terrain features (elevation and slope) be systematically integrated into a 2.5D path-planning cost function?

\item How does a slope-constrained B-spline smoothing algorithm refine
    trajectories without violating physical vehicle traversability limits?

    \item To what extent does the proposed 2.5D approach improve 3D path length, execution time, and energy efficiency compared to conventional 2D and full 3D planners across diverse terrain profiles?
    \item What is the sensitivity and impact of key 2.5D planning parameters (e.g., slope weight $\lambda$) on trajectory optimality and computational overhead?
\end{enumerate}

\section{Scope of the Study}

To ensure a focused and rigorous validation of the proposed navigation framework, the scope of this study is systematically bounded across the following operational and architectural parameters:
\begin{enumerate}
    \item Environmental Mapping \& Grid Representation: The environmental mapping pipeline is restricted to converting simulation-derived 3D point cloud data (.pcd format) into static, discrete 2.5D elevation grids equipped with a single-layer vegetation penalty mask. The study does not cover real-time online mapping, dynamic obstacle perception, or full 3D voxel grid representation.
    \item Core Pathfinding Mechanics: Global path generation is limited to an 8-connected grid $A^*$ algorithm that directly evaluates 3D step distances (3DSD), localized slope magnitudes (SM), and a manually adjusted slope weight ($\lambda$) within its node-expansion cost function. Alternative search paradigms—such as continuous sampling-based planners (e.g., RRT*), dynamic local re-planning, or kinodynamic constraints during the initial graph search—fall outside the scope of this work.
    \item Trajectory Smoothing \& Safety Enforcement: Post-processing path refinement is strictly bounded to an open-uniform B-spline module, coupled with discrete sampling checks to enforce physical traversability boundaries, vegetation limits, and maximum slope thresholds. Complex vehicle dynamics, such as wheel-terrain friction, suspension deflection, or roll-stability modeling, are not incorporated into the trajectory optimization layer.
    \item Experimental Variables \& Map Configurations: System performance is evaluated across three synthetic terrain profiles—low-gradient, medium-relief, and high-relief topographies—using five specific slope weight ($\lambda$) configurations and three fixed start-goal coordinate pairs per environment. The study excludes testing on unstructured real-world agricultural fields or hardware-in-the-loop (HIL) setups.
    \item Simulation \& Validation Environment: All experimental trials, algorithmic profiling, and trajectory visualizations are conducted entirely within a virtual simulation sandbox to safely analyze performance. Consequently, physical vehicle deployment, field trials, and edge-hardware profiling on physical ground platforms are excluded from this study.
\end{enumerate}

Finally, this research delivers critical methodological advancements by
establishing a highly systematic and cost-effective validation pipeline for
autonomous vehicle benchmarking. The simulation-only approach allows for
extensive, repeatable experimentation across a broad spectrum of heterogeneous
synthetic terrain profiles---specifically low-gradient plains, medium-relief
hills, and high-relief escarpments---completely bypassing the severe logistical
constraints, equipment costs, and physical hazards of real-world prototyping
trials. Utilizing a comprehensive performance suite comprising exact 3D path
lengths ($L_{3D}$), travel-time proxies (TTP), computational execution runtimes,
and energy consumption proxies (ECP) provides a rigorous mathematical baseline
for direct comparison against traditional 2D planners. Benchmarking these
metrics across a variable range of slope weight configurations offers
transparent, quantifiable insights into the precise efficiency and safety gains
achieved by the terrain-aware approach. This methodological rigor ensures that
the resulting data trends remain robustly scalable and mathematically sound for
future physical implementations.
